\medskip

%Dans le cadre d’un projet de labellisation \og Éducation au développement durable \fg, un collège
%réalise deux enquêtes sur une période donnée.
%
%\medskip

\begin{enumerate}
\item %La première enquête porte sur le gaspillage alimentaire à la cantine.

%Pendant sept semaines, on relève la masse totale, en kilogramme, d'aliments jetés chaque
%semaine :
%
%\begin{center}
%\begin{tabularx}{\linewidth}{|l|*{7}{>{\centering \arraybackslash}X|}}\hline
%Semaine &	1 &2 &3 &4 &5 &6 &7\\ \hline
%Masse (kg) &62 &59 &74 &68 &55 &61 &71\\ \hline
%\end{tabularx}
%\end{center}
%
%Ce collège s'est donné comme objectif que la moyenne, par semaine, de déchets
%alimentaires sur les 7 semaines ne dépasse pas 65~kg.
%
%Montrer que ce collège a atteint son objectif.

Moyenne des déchets sur la semaine : 

$\dfrac{62 + 59 + + 74 + 68 + 55 + 61 + 71}{7} = \dfrac{450}{7} \approx 64,3 < 65.$
\item %La seconde enquête porte sur les déplacements des élèves à vélo entre le domicile et le
%collège.

%Le diagramme ci-dessous représente, pour chaque distance, l'effectif des élèves qui
%parcourent cette distance en vélo pour aller au collège. (Les élèves qui n’utilisent pas le vélo
%pour se rendre au collège parcourent 0 km à vélo.)
%
%\begin{center}
%\psset{xunit=1cm,yunit=0.15cm}
%\begin{pspicture}(-1.6,-3)(9.3,50)
%\rput(4,48){\textbf{Distance parcourue à vélo par les élèves du collège}}
%\uput[d](4,-2){\small Distance parcourue en km}
%\rput{90}(-1,23){\small Effectif}
%\psaxes[linewidth=1pt,Dx=10,Dy=2,labelFontSize=\scriptstyle](0,0)(0,0)(9.3,46)
%\multido{\n=0+1,\na=1+1}{9}{\uput[d](\na,0){\footnotesize \n}}
%\psframe[fillstyle=solid,fillcolor=purple](0.7,0)(1.3,33)
%\psframe[fillstyle=solid,fillcolor=purple](1.7,0)(2.3,32)
%\psframe[fillstyle=solid,fillcolor=purple](2.7,0)(3.3,42)
%\psframe[fillstyle=solid,fillcolor=purple](3.7,0)(4.3,31)
%\psframe[fillstyle=solid,fillcolor=purple](4.7,0)(5.3,36)
%\psframe[fillstyle=solid,fillcolor=purple](5.7,0)(6.3,28)
%\psframe[fillstyle=solid,fillcolor=purple](6.7,0)(7.3,24)
%\psframe[fillstyle=solid,fillcolor=purple](7.7,0)(8.3,22)
%\psframe[fillstyle=solid,fillcolor=purple](8.7,0)(9.3,14)
%\end{pspicture}
%\end{center}

	\begin{enumerate}
	\item %Déterminer l'effectif total d'élèves de ce collège.
	L'effectif total du collège est égal à :
	
	$33 + 32 + 42 + + 30 + 36 +  + 28 +  24 + 22 + 14 = 243$.
	\item %Pour ce collège, l'affirmation \og Plus de 30\,\% des élèves ont parcouru au moins 5~km à vélo
%pour se rendre au collège \fg est-elle vraie ?

Ont parcouru 5 km ou plus à vélo : 

$28 + 24 + 22 + 14 = 88$ élèves.

D'autre part 30\,\% des 243 élèves représentent : $\dfrac{30}{100} \times 243 = 72,3$.

Le premier entier supérieur à ce nombre est 73,et comme $88 > 73$, l'affirmation est vraie.

%\textbf{Justifier la réponse en précisant la démarche.}
	\end{enumerate}
\end{enumerate}

\bigskip
